% संपूर्ण मराठी ग्रंथातील प्रकरण 25: ट्यूरिंग यंत्रांवरील संगणने
% हा संपादनयोग्य स्रोतखंड आहे. संकलनासाठी संपूर्ण स्रोतसंग्रहातील
% अक्षररूपे, पूर्वभाग, चिन्हव्याख्या आणि आकृती-स्रोत आवश्यक आहेत.
% मूळ: Open Logic Project; मजकूर आणि रूपांतर CC BY 4.0.
% यंत्रानुवाद, दुरुस्ती आणि तपासणी: OpenAI Codex —
% GPT-5.6 Sol आणि GPT-6 Sol, दोन्ही Ultra effort.
% दुरुस्ती-पुनरावलोकन: GPT-6.1 Sol, Ultra effort.
\chapter{ट्यूरिंग यंत्रांवरील संगणने}\label{tur:mac::chap}








\section{परिचय}\label{tur:mac:int:sec}

$\Nat$ पासून $\Nat$ मध्ये जाणारे एखादे फलन \emph{संगणनक्षम} असणे म्हणजे
नेमके काय? या प्रश्नाला सर्वप्रथम दिलेल्या आणि सर्वाधिक परिचित
उत्तरांपैकी एक असे आहे: एखादे फलन ट्यूरिंग यंत्राने काढता येत असेल, तर ते
संगणनक्षम असते. ॲलन ट्यूरिंग यांनी 1936 मध्ये ही संकल्पना मांडली.
ट्यूरिंग यंत्र हे \emph{संगणनाचे प्रतिमान} आहे---म्हणजे ``संगणनाची
कार्यपद्धती'' ही कल्पना गणिती नेमकेपणाने परिभाषित करण्याचा एक मार्ग.
त्या कल्पनेचा नेमका अर्थ काय, यावर मतभेद असले तरी ट्यूरिंग यंत्रे ही
संगणनाच्या कार्यपद्धती निर्दिष्ट करण्याची एक पद्धत आहेत, यावर व्यापक
सहमती आहे. ``ट्यूरिंग यंत्र'' या नावामुळे हलणारे भाग असलेल्या भौतिक
यंत्राची प्रतिमा डोळ्यांसमोर येऊ शकते; पण काटेकोरपणे पाहता ट्यूरिंग
यंत्र ही पूर्णपणे गणिती रचना आहे. म्हणून ती संगणनाच्या कार्यपद्धतीची
आदर्शीकृत मांडणी करते. उदाहरणार्थ, ट्यूरिंग यंत्राला लागणाऱ्या वेळेवर
किंवा स्मृतीवर आपण कोणतेही बंधन घालत नाही: एखाद्या संगणनाला विश्वातील
अणूंपेक्षा जास्त साठवण-जागा किंवा पायऱ्या लागणार असल्या तरी ट्यूरिंग
यंत्र ते संगणन करू शकते.

\begin{explain}
ट्यूरिंग यंत्राचा विचार एका विशिष्ट प्रकारच्या काल्पनिक यंत्रणेसाठीचा
कार्यक्रम म्हणून करणे बहुधा सर्वांत सोयीचे ठरेल. या यंत्रणेत एक
\emph{फीत} आणि एक \emph{वाचन-लेखन शीर्ष} असते. आपल्या ट्यूरिंग
यंत्रांच्या आवृत्तीत फीत एका दिशेने (उजवीकडे) अनंत असते आणि ती
\emph{घरांमध्ये} विभागलेली असते. प्रत्येक घरात सांत \emph{वर्णमालेतील}
एखादे चिन्ह असू शकते. अशा वर्णमालेत कितीही भिन्न चिन्हे असू शकतात;
पण आपल्याला मुख्यतः तीनच पुरेशी पडतील: $\TMendtape$, $\TMblank$ आणि
$\TMstroke$. यंत्रणा सुरू करताना फीत रिकामी असते (म्हणजे प्रत्येक घरात
$\TMblank$ हे चिन्ह असते); अपवाद म्हणजे सर्वांत डावीकडील घर, ज्यात
$\TMendtape$ असते, आणि \emph{आदान} असलेली सांत संख्येने घरे. कोणत्याही
क्षणी यंत्रणा सांत संख्येतील एखाद्या \emph{अवस्थेत} असते. सुरुवातीला
शीर्ष आदानाच्या पहिल्या घरावर, म्हणजे फितीच्या डाव्या टोकाच्या
घराच्या लगेच उजवीकडील घरावर, वाचत असते आणि यंत्रणा निर्दिष्ट
\emph{प्रारंभिक अवस्थेत} असते. यंत्रणेच्या प्रत्येक पायरीवर शीर्ष
सध्या वाचत असलेल्या घरातील चिन्ह, यंत्रणेची अवस्था आणि ट्यूरिंग
यंत्राचा कार्यक्रम मिळून पुढे काय होईल ते ठरवतात. हा कार्यक्रम एका
आंशिक फलनाने दिलेला असतो: त्याला अवस्था~$q$ आणि चिन्ह~$\sigma$ ही
आदाने दिल्यावर ते $\tuple{q', \sigma',
  D}$ ही त्रयी देते. यंत्रणा $q$ अवस्थेत असताना $\sigma$ चिन्ह
वाचत असेल, तर ती सध्याच्या घरातील चिन्हाच्या जागी $\sigma'$ लिहिते;
$D$ हे $\TMleft$, $\TMright$ किंवा $\TMstay$ यांपैकी कोणते आहे,
त्यानुसार शीर्ष डावीकडे, उजवीकडे किंवा त्याच जागी राहते; आणि यंत्रणा
$q'$ अवस्थेत जाते.


उदाहरणार्थ, \ref{tur:mac:int:fig:tm} मधील स्थिती पाहा.
\begin{figure}
  \olasset{assets/diagrams/turing-machine.tikz}
  \caption{कार्यक्रम चालवत असलेले ट्यूरिंग यंत्र.}
  \label{tur:mac:int:fig:tm}
\end{figure}
ट्यूरिंग यंत्राच्या फितीच्या दिसणाऱ्या भागात सर्वांत डावीकडील घरात
फितीच्या टोकाचे $\TMendtape$ चिन्ह आहे; त्यानंतर तीन $1$, एक $0$ आणि
आणखी चार $1$ आहेत. शीर्ष डावीकडून तिसरे घर वाचत आहे; त्या घरात
$1$ आहे आणि यंत्रणा $q_1$ अवस्थेत आहे---म्हणजे ``यंत्र $q_1$
अवस्थेत $1$ वाचत आहे'' असे आपण म्हणतो. ट्यूरिंग यंत्राचा कार्यक्रम
$\tuple{q_1, 1}$ या आदानासाठी $\tuple{q_2,
0, \TMstay}$ ही त्रयी देत असेल, तर यंत्र तिसऱ्या घरातील $1$ च्या
जागी $0$ लिहील, वाचन-लेखन शीर्ष त्याच जागी ठेवील आणि $q_2$
अवस्थेत जाईल. त्यानंतर कार्यक्रम $\tuple{q_2, 0}$ या आदानासाठी
$\tuple{q_3, 0, \TMright}$ देत असेल, तर यंत्र त्या $0$ च्या
जागी पुन्हा $0$ लिहील (म्हणजे प्रत्यक्षात शीर्षाखालील फितीतील मजकूर
अपरिवर्तित राहील), एक घर उजवीकडे सरकेल आणि $q_3$ अवस्थेत जाईल.
अशाच रीतीने पुढे चालेल.

एखादी अवस्था $q_n$ आणि चिन्ह $\sigma$ अशी आढळली की
$\tuple{q_n, \sigma}$ साठी कोणतीही सूचना नाही---म्हणजे
$\tuple{q_n,\sigma}$ या आदानावर संक्रमण फलन अपरिभाषित आहे---तर यंत्र
\emph{थांबते} असे आपण म्हणतो.
दुसऱ्या शब्दांत, यंत्राला अमलात आणण्यासाठी कोणतीही सूचना उरत नाही
आणि त्या टप्प्यावर त्याचे कार्य थांबते. काही वेळा थांबणे एका विशिष्ट
थांबण्याच्या अवस्थेने~$h$ दर्शवतात. याचे अधिक तपशीलवार प्रदर्शन पुढे
केले जाईल.
\end{explain}

\begin{digress}
ट्यूरिंग यांच्या ``On computable numbers'' या शोधनिबंधाचे वैशिष्ट्य
असे की त्यात त्यांनी केवळ औपचारिक व्याख्याच दिली नाही, तर ती व्याख्या
संगणनक्षमतेची अंतर्ज्ञानी कल्पना योग्यरीत्या पकडते, असा युक्तिवादही
केला. व्याख्येवरून ट्यूरिंग यंत्राने काढता येणारे कोणतेही फलन अंतर्ज्ञानी
अर्थाने संगणनक्षम आहे, हे स्पष्ट व्हावे. ट्यूरिंग यांनी याचा उलट
दिशेचा दावा---म्हणजे आपण स्वाभाविकपणे संगणनक्षम मानू असे कोणतेही फलन
अशा यंत्राने काढता येते---यासाठी तीन प्रकारचे युक्तिवाद दिले. त्यांच्या
शब्दांत ते असे:
\begin{enumerate}
\item अंतःप्रज्ञेला थेट आवाहन.
\item दोन व्याख्या सममूल्य आहेत, याची सिद्धता (नवी व्याख्या
  अंतःप्रज्ञेला अधिक पटत असल्यास).
\item संगणनक्षम संख्यांच्या मोठ्या वर्गांची उदाहरणे देणे.
\end{enumerate}
आपला उद्देश संगणनक्षमता ``तत्त्वतः'' परिभाषित करण्याचा आहे: म्हणजे
वेळ आणि जागा यांच्या व्यावहारिक मर्यादा विचारात न घेता. अर्थात,
संगणनक्षमतेची ही सर्वांत व्यापक व्याख्या मिळाल्यानंतर मर्यादित
संसाधनांतील संगणनाचाही अभ्यास करता येतो. ``संगणकीय जटिलता'' म्हणून
ओळखल्या जाणाऱ्या विषयाचा हा मध्यवर्ती भाग आहे.
\end{digress}

\begin{history}
ॲलन ट्यूरिंग यांनी 1936 मध्ये ट्यूरिंग यंत्रांची कल्पना मांडली.
त्या काळी त्यांना प्रथम-क्रम तर्कशास्त्रातील निर्णयक्षमतेचा प्रश्न
विशेष महत्त्वाचा वाटत होता; तरी संगणकाच्या रचनेच्या पायावरील निर्णायक
शोधनिबंध म्हणून त्यांच्या त्या निबंधाचे वर्णन झाले आहे. त्या निबंधात
ट्यूरिंग यांनी संगणनक्षम वास्तव संख्यांवर---म्हणजे ज्यांचे दशांश
विस्तार संगणनक्षम आहेत अशा संख्यांवर---लक्ष केंद्रित केले. पण
नैसर्गिक संख्यांवरील संगणनक्षम फलनांसाठी आणि पुढील गोष्टींसाठीही या
संकल्पना रूपांतरित करणे कठीण नाही, असे त्यांनी नमूद केले. लक्षात घ्या:
पहिला कार्यरत सर्वसाधारण वापराचा संगणक 1941 मध्ये (जर्मन कॉनराड
झुसे यांनी बर्लिनमध्ये) बांधला जाण्याच्या पूर्ण पाच
वर्षे आधी हे झाले. 1936 नंतर सात वर्षांनी (1943) टॉमी फ्लावर्स आणि
डॉलिस हिलमधील त्यांच्या सहकाऱ्यांनी संकेतभेदनासाठी कोलॉसस बांधला;
तो ब्लेच्ली पार्कमध्ये 1944 मध्ये कार्यरत झाला.
अमेरिकन ENIAC
(1945) हा नऊ वर्षांनी आला; मँचेस्टरमधील प्रायोगिक
यंत्र---Manchester Small-Scale Experimental Machine---बारा
वर्षांनी (1948) बांधले गेले; आणि अमेरिकन BINAC (1949) ची पहिली
चाचणी तेरा वर्षांनी झाली. Manchester SSEM हा संग्रहित-कार्यक्रम
असलेला पहिला इलेक्ट्रॉनिक संगणक होता---त्यात कार्यक्रमाच्या सूचना
इलेक्ट्रॉनिक स्मृतीत ठेवता येत होत्या. त्याआधीच्या यंत्रांत सूचना
कागदी फितीवर दिल्या जात किंवा काही यंत्रांत तारा नव्याने जोडाव्या लागत.

\end{history}











\section{ट्यूरिंग यंत्रांचे निरूपण}\label{tur:mac:rep:sec}

\begin{explain}
ट्यूरिंग यंत्रे \emph{अवस्था-आलेखांद्वारे} चित्ररूपात दाखवता येतात.
या आलेखांत बाणांनी जोडलेली अवस्थांची घरे असतात. अर्थातच प्रत्येक
अवस्था-घर यंत्राची एक अवस्था दर्शवते. प्रत्येक बाण त्या अवस्थेतून
अमलात आणता येणारी एक सूचना दर्शवतो; सूचनेचे तपशील त्या बाणाच्या
वर किंवा खाली लिहिलेले असतात. पुढील यंत्र पाहा. त्याच्या $q_0$ आणि
$q_1$ अशा फक्त दोन अंतर्गत अवस्था आणि एकच सूचना आहे:
\[\begin{tikzpicture}[->,>=stealth',shorten >=1pt,auto,node distance=2.8cm,
                    semithick]
  \tikzstyle{every state}=[fill=none,draw=black,text=black]
  \node[initial,state]   (A)              {$q_0$};
  \node[state]   (B) [right of=A] {$q_1$};
  \path (A) edge  node {\TMtrans{\TMblank}{\TMstroke}{\TMright}} (B);
\end{tikzpicture}
\]
ट्यूरिंग यंत्राला वाचन-लेखन शीर्ष आणि त्यावर आदान लिहिलेली फीत
असते, हे आठवा. या सूचनेचा अर्थ असा: \emph{$q_0$ अवस्थेत
रिकामे चिन्ह~$\TMblank$ वाचत असाल, तर $\TMstroke$ लिहा,
उजवीकडे सरका आणि $q_1$ अवस्थेत जा}. हे संक्रमण फलनाने
$\tuple{q_0, \TMblank}$ याचा $\tuple{q_1, \TMstroke,
\TMright}$ शी केलेला संबंध याच्या सममूल्य आहे.
\end{explain}

\begin{ex}
\emph{सम यंत्र}: फितीवर $\TMstroke$ चिन्हांची संख्या सम असेल तेव्हा
आणि तेव्हाच पुढील ट्यूरिंग यंत्र थांबते (फितीवरील पहिल्या
$\TMblank$ आधी सर्व $\TMstroke$ चिन्हे येतात, हे गृहीत धरून).
\[\begin{tikzpicture}[->,>=stealth',shorten >=1pt,auto,node distance=2.8cm,
                    semithick]
  \tikzstyle{every state}=[fill=none,draw=black,text=black]
  \node[initial,state]         (A)              {$q_0$};
  \node[state]         (B) [right of=A] {$q_1$};
  \path (A) edge [bend left] node {\TMtrans{\TMstroke}{\TMstroke}{\TMright}} (B)
        (B) edge [loop above] node {\TMtrans{\TMblank}{\TMblank}{\TMright}} (B)
            edge [bend left] node {\TMtrans{\TMstroke}{\TMstroke}{\TMright}} (A);
\end{tikzpicture}
\]

हा अवस्था-आलेख पुढील संक्रमण फलनाशी संबंधित आहे:
\begin{align*}
\delta(q_0, \TMstroke) & = \tuple{q_1, \TMstroke, \TMright},\\
\delta(q_1, \TMstroke) & = \tuple{q_0, \TMstroke, \TMright},\\
\delta(q_1, \TMblank)  & = \tuple{q_1, \TMblank, \TMright}
\end{align*}
\end{ex}

\begin{explain}
वरील यंत्र आदानातील रेघांची संख्या सम असेल तेव्हाच थांबते.
अन्यथा यंत्र (सैद्धांतिकदृष्ट्या) अमर्याद काळ कार्यरत राहते.
कोणत्याही यंत्राच्या कोणत्याही आदानासाठी त्याच्या \emph{स्थितिवर्णनांचा}
क्रम पाहता येतो; यंत्र थांबल्यास या क्रमावरून त्याचे प्रदान निश्चित
करता येते. पण यंत्र थांबत नसेल तर स्थितिवर्णनांचा क्रम अनंत असतो;
त्यावरून प्रत्येक आदानासाठी सांत वेळेत प्रदान मिळते असा अर्थ नाही.

स्थितिवर्णनाची औपचारिक
व्याख्या पुढे देऊ. सध्या, यंत्र चालू असताना प्रत्येक क्षणी त्याची
स्थिती दाखवणाऱ्या आकृत्यांची मालिका असे स्थितिवर्णनांचे अंतर्ज्ञानी
आकलन करता येईल. स्थितिवर्णनात फितीवरील मजकूर, यंत्राची अवस्था
आणि वाचन-लेखन शीर्षाची जागा दिसते.

सम यंत्राला चार $\TMstroke$ चिन्हांचे आदान दिल्यास त्याची
स्थितिवर्णने पाहू. या वेळी यंत्र थांबेल, अशी अपेक्षा आहे. मग
तीन $\TMstroke$ चिन्हांचे आदान देऊ; त्या वेळी ते कायम चालू राहील.

यंत्र $q_0$ अवस्थेत सुरू होते आणि सर्वांत डावीकडील $\TMstroke$
वाचते. त्याची प्रारंभिक स्थिती अशी दर्शवता येते:
\[\TMendtape \TMstroke_0 \TMstroke \TMstroke \TMstroke \TMblank \ldots
\]
वरील स्थितिवर्णन सरळ आहे. दिसते त्याप्रमाणे यंत्र $q_0$ अवस्थेत
सुरू होते आणि सर्वांत डावीकडील $\TMstroke$ वाचते. पहिल्या

$\TMstroke$ च्या खाली लिहिलेला अवस्थेच्या नावाचा निर्देशांक हे
दर्शवतो. या टप्प्यावर लागू होणारी सूचना
$\delta(q_0, \TMstroke) = \tuple{q_1, \TMstroke, \TMright}$ आहे;
म्हणून यंत्र फितीवर उजवीकडे सरकते आणि $q_1$ अवस्थेत जाते.
\[\TMendtape \TMstroke \TMstroke_1 \TMstroke \TMstroke \TMblank \ldots
\]
आता यंत्र $q_1$ अवस्थेत $\TMstroke$ वाचत आहे, म्हणून
$\delta(q_1, \TMstroke) = \tuple{q_0,
  \TMstroke, \TMright}$ ही सूचना ``पाळावी'' लागते. त्यातून पुढील
स्थितिवर्णन मिळते:
\[\TMendtape \TMstroke \TMstroke \TMstroke_0 \TMstroke \TMblank \ldots
\]
यंत्र पुढे चालू राहिल्यावर हे नियम त्याच क्रमाने पुन्हा लागू होतात
आणि पुढील दोन स्थितिवर्णने मिळतात:
\[\TMendtape \TMstroke \TMstroke \TMstroke \TMstroke_1 \TMblank \ldots
\]
\[\TMendtape \TMstroke \TMstroke \TMstroke \TMstroke \TMblank_0 \ldots
\]
आता यंत्र $q_0$ अवस्थेत $\TMblank$ वाचत आहे. संक्रमण-आलेखावरून
येथे अमलात आणण्यासाठी कोणतीही सूचना नाही, हे सहज दिसते;
म्हणून यंत्र थांबले आहे. याचा अर्थ आदान स्वीकारले गेले.

आता यंत्राला तीन $\TMstroke$ चिन्हांचे आदान देऊन सुरू करू.
सुरुवातीची काही स्थितिवर्णने मिळतीजुळती असतील, कारण त्याच सूचना
अमलात येतात; फक्त फितीवरील आदान थोडे वेगळे आहे:
\[\TMendtape \TMstroke_0 \TMstroke \TMstroke \TMblank \ldots
\]
\[\TMendtape \TMstroke \TMstroke_1 \TMstroke \TMblank \ldots
\]
\[\TMendtape \TMstroke \TMstroke \TMstroke_0 \TMblank \ldots
\]
\[\TMendtape \TMstroke \TMstroke \TMstroke \TMblank_1 \ldots
\]
आता यंत्र सर्व $\TMstroke$ चिन्हांच्या पुढे गेले आहे आणि $q_1$
अवस्थेत $\TMblank$ वाचत आहे. आलेखात दाखवल्याप्रमाणे
$\delta(q_1, \TMblank) =\tuple{q_1,
\TMblank, \TMright}$ या रूपाची सूचना आहे. फीत उजवीकडे अनंतपर्यंत
$\TMblank$ चिन्हांनी भरलेली असल्याने यंत्र ही सूचना \emph{कायम}
अमलात आणत राहील; ते $q_1$ अवस्थेतच राहून सतत उजवीकडे सरकेल.
यंत्र कधीच थांबणार नाही आणि आदान स्वीकारणार नाही.
\end{explain}

\begin{explain}
सर्व यंत्रे थांबतातच असे नाही, हे लक्षात ठेवणे महत्त्वाचे आहे. थांबणे
म्हणजे अमलात आणण्यासाठी सूचना संपणे असेल, तर प्रत्येक अवस्थेत
प्रत्येक चिन्हासाठी बाहेर जाणारा बाण ठेवून कधीही न थांबणारे यंत्र
बनवता येते. $q_0$ अवस्थेत $\TMblank$ वाचण्यासाठी एक सूचना
जोडून सम यंत्र कायम चालू राहील असे बदलता येते.
\end{explain}

\begin{ex}
\[\begin{tikzpicture}[->,>=stealth',shorten >=1pt,auto,node distance=2.8cm,
                    semithick]
  \tikzstyle{every state}=[fill=none,draw=black,text=black]
  \node[initial,state] (A)              {$q_0$};
  \node[state]         (B) [right of=A] {$q_1$};
  \path
  (A) edge [bend left] node {\TMtrans{\TMstroke}{\TMstroke}{\TMright}} (B)
      edge [loop above] node {\TMtrans{\TMblank}{\TMblank}{\TMright}} (A)
  (B) edge [loop above] node {\TMtrans{\TMblank}{\TMblank}{\TMright}} (B)
      edge [bend left] node {\TMtrans{\TMstroke}{\TMstroke}{\TMright}} (A);
\end{tikzpicture}
\]
\end{ex}

\begin{explain}
यंत्र-सारणी हा ट्यूरिंग यंत्रांचे निरूपण करण्याचा दुसरा मार्ग आहे.
यंत्र-सारणीत फितीवरील वर्णमाला $x$-अक्षावर आणि यंत्राच्या अवस्था
$y$-अक्षावर दाखवलेल्या असतात. सारणीतील प्रत्येक अवस्था आणि चिन्ह
यांच्या छेदावर सूचनेचा उरलेला भाग---नवी अवस्था, नवे चिन्ह आणि
हालचालीची दिशा---लिहिलेला असतो. यंत्र कोणत्या अवस्थेत आणि कोणते
चिन्ह वाचताना थांबते हे यंत्र-सारणीवरून सहज ठरवता येते. सारणीत
जिथे नोंद नाही, तिथे यंत्र थांबू शकते. अवस्था-आलेख किंवा सूचनासंच
यांत यंत्र नेमके कुठे थांबेल हे लगेच स्पष्ट होत नाही; पण
यंत्र-सारणीत रिकामी घरे शोधून अशी ठिकाणे पटकन ओळखता येतात.
\end{explain}

\begin{ex}
सम यंत्राची यंत्र-सारणी अशी आहे:
\[\centering
\begin{tabular}{lllll}
\cline{2-4}
\multicolumn{1}{l|}{}      & \multicolumn{1}{c|}{$\TMblank$}
& \multicolumn{1}{c|}{$\TMstroke$}     & \multicolumn{1}{c|}{$\TMendtape$}   &  \\ \cline{1-4}
\multicolumn{1}{|l|}{$q_0$} & \multicolumn{1}{l|}{}
& \multicolumn{1}{l|}{\TMtrans{\TMstroke}{q_1}{\TMright}} &
\multicolumn{1}{l|}{\phantom{\TMtrans{\TMstroke}{q_1}{\TMright}}} &  \\ \cline{1-4}
\multicolumn{1}{|l|}{$q_1$} & \multicolumn{1}{l|}{\TMtrans{\TMblank}{q_1}{\TMright}}
& \multicolumn{1}{l|}{\TMtrans{\TMstroke}{q_0}{\TMright}} &
\multicolumn{1}{l|}{\phantom{\TMtrans{\TMstroke}{q_1}{\TMright}}} &  \\ \cline{1-4}
\end{tabular}
\]
दिसते त्याप्रमाणे यंत्र $q_0$ अवस्थेत $\TMblank$ वाचताना थांबते.
\end{ex}

\begin{explain}
आतापर्यंत आपण फक्त आदान वाचून ते स्वीकारणारी यंत्रे पाहिली.
परंतु ट्यूरिंग यंत्रे वाचूही शकतात आणि लिहूही शकतात. अशा यंत्राचे
एक उदाहरण (अशी अनेक उदाहरणे असली तरी) म्हणजे \emph{दुप्पट करणारे
यंत्र}. हे यंत्र फितीवरील $n$ $\TMstroke$ चिन्हांचा सलग गट
आदान म्हणून घेते आणि $2n$ $\TMstroke$ चिन्हांचा सलग गट प्रदान करते.
\end{explain}

\begin{ex}
  \label{tur:mac:rep:ex:doubler} दुप्पट करणारे यंत्र बनवण्यापूर्वी ही समस्या
  सोडवण्यासाठी \emph{धोरण} ठरवणे महत्त्वाचे आहे. यंत्राला (आपल्या
  व्याख्येनुसार) त्याने किती $\TMstroke$ चिन्हे वाचली हे लक्षात
  ठेवता येत नाही; म्हणून फितीवरील सर्व $\TMstroke$ चिन्हांचा मागोवा
  घेण्याचा मार्ग हवा. एक मार्ग म्हणजे आदान आणि प्रदान यांमध्ये
  $\TMblank$ ठेवून त्यांना वेगळे करणे. मग यंत्र आदानातील पहिले
  $\TMstroke$ पुसू शकते, आदानाचा उरलेला भाग पार करू शकते,
  $\TMblank$ सोडू शकते आणि दोन नवी $\TMstroke$ चिन्हे लिहू शकते.
  नंतर ते आदानातील दुसरे $\TMstroke$ शोधून त्याच्यासाठीही असेच
  दुपटीकरण करील. आदानातील प्रत्येक एका $\TMstroke$ साठी ते
  प्रदानात दोन $\TMstroke$ चिन्हे लिहील. पुढे जाताना आदान
  पुसल्यामुळे कोणतेही $\TMstroke$ चिन्ह सुटणार नाही किंवा त्याची
  दुप्पट मोजणी होणार नाही, याची खात्री होते. संपूर्ण आदान पुसल्यावर
  फितीवर $2n$ $\TMstroke$ चिन्हे उरतील. परिणामी ट्यूरिंग यंत्राचा
  अवस्था-आलेख \ref{tur:mac:rep:fig:doubler} मध्ये दाखवला आहे.
  \begin{figure}
\[\begin{tikzpicture}[->,>=stealth',shorten >=1pt,auto,node distance=2.8cm,
                    semithick]
  \tikzstyle{every state}=[fill=none,draw=black,text=black]
  \node[initial,state] (1)              {$q_0$};
  \node[state]         (2) [right of=1] {$q_1$};
  \node[state]         (3) [right of=2] {$q_2$};
  \node[state]         (4) [below of=3] {$q_3$};
  \node[state]         (5) [left of=4]  {$q_4$};
  \node[state]         (6) [left of=5]  {$q_5$};
  \path (1) edge node {\TMtrans{\TMstroke}{\TMblank}{\TMright}} (2)
    (2) edge [loop above] node {\TMtrans{\TMstroke}{\TMstroke}{\TMright}} (2)
      edge node {\TMtrans{\TMblank}{\TMblank}{\TMright}} (3)
    (3) edge [loop above] node {\TMtrans{\TMstroke}{\TMstroke}{\TMright}} (3)
        edge node {\TMtrans{\TMblank}{\TMstroke}{\TMright}} (4)
    (4) edge [loop below] node {\TMtrans{\TMblank}{\TMstroke}{\TMleft}} (4)
        edge node {\TMtrans{\TMstroke}{\TMstroke}{\TMleft}} (5)
    (5) edge [loop below]  node {\TMtrans{\TMstroke}{\TMstroke}{\TMleft}} (5)
        edge              node {\TMtrans{\TMblank}{\TMblank}{\TMleft}} (6)
    (6) edge [loop below] node {\TMtrans{\TMstroke}{\TMstroke}{\TMleft}} (6)
        edge              node {\TMtrans{\TMblank}{\TMblank}{\TMright}} (1);
\end{tikzpicture}
\]
\caption{दुप्पट करणारे यंत्र}
\label{tur:mac:rep:fig:doubler}
\end{figure}
\end{ex}

\begin{prob}
एखादे आदान निवडा आणि \ref{tur:mac:rep:ex:doubler} मधील दुप्पट
करणाऱ्या यंत्राची स्थितिवर्णने क्रमाने शोधा.
\end{prob}

\begin{prob}
$\{\TMendtape,\TMblank, A, B\}$ ही वर्णमाला असलेले ट्यूरिंग यंत्र
रचा. $A$ आणि $B$ यांच्या कोणत्याही चिन्हमालेत $A$ ची संख्या $B$
च्या संख्येइतकी \emph{आणि} सर्व $A$ सर्व $B$ च्या आधी असतील,
तर यंत्राने ती चिन्हमाला स्वीकारावी, म्हणजे तिच्यावर थांबावे.
$A$ आणि $B$ यांच्या संख्या असमान असतील किंवा सर्व $A$ सर्व $B$
च्या आधी नसतील, तर यंत्राने ती नाकारावी, म्हणजे तिच्यावर थांबू नये.
(उदाहरणार्थ, यंत्राने $AABB$ आणि $AAABBB$ स्वीकाराव्यात; पण
$AAB$ आणि $AABBAABB$ या दोन्ही नाकाराव्यात.)
\end{prob}

\begin{prob}
$\{\TMendtape,\TMblank, A, B\}$ ही वर्णमाला असलेले ट्यूरिंग यंत्र
रचा. त्याने $A$ आणि $B$ यांची कोणतीही चिन्हमाला $\alpha$ आदान
म्हणून घेऊन तिची प्रत तयार करावी आणि $\alpha\alpha$ या रूपाचे
प्रदान द्यावे. (उदाहरणार्थ, $ABBA$ या आदानापासून $ABBAABBA$
हे प्रदान मिळावे.)
\end{prob}

\begin{prob}
\emph{वर्णक्रमात आहे का?}: $\{\TMendtape,\TMblank, A, B\}$ ही
वर्णमाला असलेले ट्यूरिंग यंत्र रचा. $A$ आणि $B$ यांची सांत
क्रमिका आदान दिल्यावर सर्व $A$ सर्व $B$ यांच्या डावीकडे आहेत की
नाही, हे त्याने तपासावे. यंत्राने आदान-चिन्हमाला फितीवर ठेवावी;
ती ``वर्णक्रमात'' असेल तर थांबावे आणि नसेल तर कायम पुनरावर्तनात
चालू राहावे.
\end{prob}

\begin{prob}
\emph{वर्णक्रम लावणारे यंत्र}: $\{\TMendtape,\TMblank, A, B\}$
ही वर्णमाला असलेले ट्यूरिंग यंत्र रचा. $A$ आणि $B$ यांची सांत
क्रमिका आदान म्हणून घेऊन त्यांची पुनर्रचना अशा रीतीने करावी की
सर्व $A$ सर्व $B$ यांच्या डावीकडे येतील. (उदाहरणार्थ, $BABAA$
ही क्रमिका $AAABB$ व्हावी आणि $ABBABB$ ही क्रमिका $AABBBB$
व्हावी.)
\end{prob}











\section{ट्यूरिंग यंत्रे}\label{tur:mac:tur:sec}

\begin{explain}
ट्यूरिंग यंत्राची औपचारिक व्याख्या अमूर्त वाटते, पण प्रत्यक्षात
ती सोपी आहे: यंत्राचे कार्य कसे चालते हे निर्दिष्ट करण्यासाठी लागणारी
सर्व माहिती ती एका गणिती रचनेत एकत्र करते. यात (1) यंत्र कोणत्या
अवस्थांत असू शकते, (2) फितीवर कोणती चिन्हे असू शकतात, (3) यंत्राने
कोणत्या अवस्थेत सुरू व्हावे आणि (4) त्याचा सूचनासंच कोणता आहे,
या गोष्टी येतात.
\end{explain}

\begin{defn}[ट्यूरिंग यंत्र]
\emph{ट्यूरिंग यंत्र} $M$ म्हणजे पुढील घटक असलेली
$\langle Q, \Sigma, q_0,
\delta\rangle$ ही क्रमित चौकडी:
\begin{enumerate}
\item \emph{अवस्थांचा} सांत संच~$Q$,
\item $\TMendtape$ आणि $\TMblank$ यांचा समावेश असलेली सांत
  \emph{वर्णमाला} $\Sigma$,
\item \emph{प्रारंभिक अवस्था}~$q_0 \in Q$,
\item सांत \emph{सूचनासंच}~$\delta\colon Q \times \Sigma
  \pto Q \times \Sigma \times \{\TMleft, \TMright, \TMstay\}$.
\end{enumerate}
या आंशिक फलनाला~$\delta$ हे $M$ चे \emph{संक्रमण फलन} असेही म्हणतात.
\end{defn}

\begin{explain}
फीत फक्त एकाच दिशेने अनंत आहे, असे आपण मानतो. म्हणून फितीचे
डावे टोक दर्शवण्यासाठी $\TMendtape$ हे खास चिन्ह ठेवणे उपयुक्त ठरते.
त्यामुळे ट्यूरिंग यंत्राच्या कार्यक्रमाला फितीच्या बाहेर जाण्याचा
``धोका'' कधी आहे हे ओळखणे सोपे जाते. हे चिन्ह कधीही बदलून लिहिले
जात नाही, असे आपण गृहीत धरू शकलो असतो; म्हणजे
$\delta(q,\TMendtape)$ परिभाषित असेल, तर
$\delta(q,\TMendtape) = \tuple{q', \TMendtape, x}$ असेल.
काही पाठ्यपुस्तके असे करतात; आपण मात्र हे गृहीत धरत नाही. ट्यूरिंग
यंत्र बनवताना ते $\TMendtape$ कधीही बदलून लिहिणार नाही याची
सावधगिरी तुम्ही घेऊ शकता. शिवाय काही प्रसंगी असे बदलून लिहिण्यास
परवानगी दिल्यास सोयीची लवचीकता मिळते.
\end{explain}

\begin{ex}
\emph{सम यंत्र:} सम यंत्र औपचारिकरीत्या
$\tuple{Q, \Sigma, q_0, \delta}$ ही क्रमित चौकडी आहे. येथे
\begin{align*}
Q & = \{ q_0, q_1 \} \\
\Sigma & = \{ \TMendtape, \TMblank, \TMstroke \}, \\
\delta(q_0, \TMstroke) & = \tuple{q_1, \TMstroke, \TMright},\\
\delta(q_1, \TMstroke) & = \tuple{q_0, \TMstroke, \TMright},\\
\delta(q_1, \TMblank)  & = \tuple{q_1, \TMblank, \TMright}.
\end{align*}
\end{ex}











\section{स्थितिवर्णने आणि संगणने}\label{cmp:tur:con:sec}

\begin{explain}
मागे सम यंत्राच्या स्थितिवर्णनांचा क्रम कसा पाहिला, ते आठवा.
फीत, वाचन-लेखन शीर्ष आणि ट्यूरिंग यंत्राचा कार्यक्रम यांची बनलेली
काल्पनिक यंत्रणा ही ट्यूरिंग यंत्राचे संगणन कसे असते ते डोळ्यांसमोर
आणण्याचा केवळ अंतर्ज्ञानी मार्ग आहे. औपचारिकरीत्या दिलेल्या आदानावर
ट्यूरिंग यंत्राचे संगणन \emph{स्थितिवर्णनांची} क्रमिका म्हणून
परिभाषित करता येते. प्रत्येक स्थितिवर्णनात चिन्हांची एक क्रमिका
(संगणनाच्या त्या टप्प्यावर फितीतील मजकुराशी संबंधित), वाचन-लेखन
शीर्षाची जागा दर्शवणारी संख्या आणि एक अवस्था असते. त्यांच्या
आधारे ट्यूरिंग यंत्र~$M$ दिलेल्या आदानावर काय काढते हे परिभाषित
करता येते.
\end{explain}

\begin{defn}[स्थितिवर्णन]
ट्यूरिंग यंत्र $M = \tuple{Q, \Sigma, q_0,
\delta}$ चे \emph{स्थितिवर्णन} म्हणजे $\tuple{C, m, q}$ ही त्रयी;
येथे
\begin{enumerate}
\item $C \in \Sigma^*$ ही $\Sigma$ मधील चिन्हांची सांत क्रमिका आहे,
\item $m \in \Nat$ ही $< \len{C}$ असलेली संख्या आहे, आणि
\item $q \in Q$ आहे.
\end{enumerate}
अंतर्ज्ञानी अर्थाने क्रमिका~$C$ म्हणजे फितीतील मजकूर (सर्वांत
डावीकडील घरापासून शेवटच्या रिकाम्या नसलेल्या किंवा आधी भेट दिलेल्या
घरापर्यंतच्या सर्व घरांतील चिन्हे); $m$~म्हणजे वाचन-लेखन शीर्ष ज्या
घरावर आहे त्या घराचा क्रमांक (सर्वांत डावीकडील घराचा क्रमांक $0$
धरून); आणि $q$ म्हणजे यंत्राची सध्याची अवस्था.
\end{defn}

\begin{explain}
ट्यूरिंग यंत्राचे संभाव्य आदान ही चिन्हांची एक क्रमिका असते; बहुतेक
वेळा ती एखाद्या संख्येला काही रूपात सांकेतिक करणारी क्रमिका असते.
ट्यूरिंग यंत्राचे प्रारंभिक स्थितिवर्णन म्हणजे त्या आदानावर यंत्र
चालवायला सुरू करतानाचे स्थितिवर्णन: फितीत डाव्या टोकाचे चिन्ह असते
आणि त्याच्या लगेच उजवीकडील घरांत आदान लिहिलेले असते; वाचन-लेखन
शीर्ष आदानाच्या सर्वांत डावीकडील घरावर असते (म्हणजे डाव्या टोकाच्या
चिन्हाच्या उजवीकडील घरावर); आणि यंत्रणा निर्दिष्ट प्रारंभिक
अवस्थेत~$q_0$ असते.

\end{explain}

\begin{defn}[प्रारंभिक स्थितिवर्णन]
आदान रिकामे नसताना, $I \in \Sigma^*$ या आदानासाठी $M$ चे \emph{प्रारंभिक स्थितिवर्णन}
असे आहे:
\[\tuple{\TMendtape \frown I, 1, q_0}.
\]
रिकामे आदान घेतल्यास फितीवरील पहिल्या रिकाम्या घराचाही
स्थितिवर्णनात समावेश करावा; अन्यथा शीर्षाचा क्रमांक आणि क्रमिकेची
लांबी यांच्यावरील वरची अट पूर्ण होत नाही.

\end{defn}

\begin{explain}
$\frown$ हे चिन्ह \emph{जोडणीसाठी} आहे---आदान-चिन्हमाला डाव्या
टोकाच्या चिन्हाच्या लगेच उजवीकडे सुरू होते.

\end{explain}

\begin{defn}
स्थितिवर्णन $\tuple{C, m, q}$ पासून \emph{एका पायरीत
$\tuple{C', m', q'}$ हे स्थितिवर्णन मिळते} ($M$ नुसार), असे आपण
तेव्हा आणि केवळ तेव्हाच म्हणतो, जेव्हा
\begin{enumerate}
\item $C$ चे $m$-वे चिन्ह $\sigma$ असते,
\item $M$ च्या सूचनासंचात $\delta(q, \sigma) =
  \tuple{q', \sigma', D}$ निर्दिष्ट असते,
\item $C'$ चे $m$-वे चिन्ह $\sigma'$ असते, आणि
\item
\begin{enumerate}
\item $D = L$ असते; या स्थितीत $m>0$ असल्यास $m' = m - 1$ असते, अन्यथा $m'=0$ असते; किंवा
\item $D = R$ आणि $m' = m + 1$ असते, किंवा
\item $D = N$ आणि $m' = m$ असते,
\end{enumerate}
\item $m' = \len{C}$ असेल, तर $\len{C'} = \len{C} + 1$ असते आणि
  $C'$ चे $m'$-वे चिन्ह~$\TMblank$ असते. अन्यथा $\len{C'}=\len{C}$ असते.
\item $i < \len{C}$ आणि $i \neq m$ असलेल्या सर्व $i$ साठी $C'(i) = C(i)$ असते,
\end{enumerate}
\end{defn}

\begin{defn}\label{cmp:tur:con:defn:run-output}
\emph{आदान~$I$ वरील $M$ ची धाव} म्हणजे $M$ च्या स्थितिवर्णनांची
$C_i$ ही सांत किंवा अनंत क्रमिका; येथे $C_0$ हे आदान~$I$ साठी $M$ चे प्रारंभिक
स्थितिवर्णन असते आणि क्रमिकेतील अंतिम स्थितिवर्णन वगळता प्रत्येक $C_i$ पासून एका पायरीत $C_{i+1}$
मिळते.


$C_k = \tuple{C, m, q}$ असेल, $C$ चे $m$-वे चिन्ह~$\sigma$ असेल
आणि $\delta(q,
\sigma)$ अपरिभाषित असेल, तर $M$ \emph{आदान $I$ वर $k$
पायऱ्यांनंतर थांबते} असे म्हणतो. त्या वेळी आदान~$I$ साठी $M$ चे
\emph{प्रदान}~$O$ असते; येथे $O$ ही शेवटी~$\TMblank$ नसलेली
चिन्हमाला असून एखाद्या~$j \in \Nat$ साठी
$C = \TMendtape \concat O \concat \TMblank^j$ असते.
($\TMblank^j$ म्हणजे $j$ $\TMblank$ चिन्हांची क्रमिका.)
ही प्रदानाची अट डाव्या टोकाचे चिन्ह शेवटपर्यंत शाबूत राहिलेल्या
धावांसाठीच थेट लागू होते. ते चिन्ह बदलून लिहिले असल्यास प्रदान
ठरवण्यासाठी अतिरिक्त नियम लागेल.

\end{defn}

\begin{explain}
या व्याख्येनुसार $M$ चे प्रदान~$O$ नेहमी~$\TMblank$ व्यतिरिक्त
इतर चिन्हाने संपते; किंवा $k$~व्या वेळी संपूर्ण फीत (सर्वांत
डावीकडील~$\TMendtape$ सोडून)~$\TMblank$ चिन्हांनी भरलेली असेल,
तर $O$~ही रिकामी चिन्हमाला असते.
\end{explain}











\section{संख्यांचे एकचिन्ही निरूपण}\label{tur:mac:una:sec}

\begin{explain}
ट्यूरिंग यंत्रे त्यांच्या फितीवर लिहिलेल्या चिन्हांच्या क्रमिकांवर
कार्य करतात. यंत्र वापरत असलेल्या वर्णमालेनुसार या क्रमिका विविध
आदानांचे आणि प्रदानांचे निरूपण करू शकतात. विशेष रस आहे तो
\emph{अंकगणितीय} फलने, म्हणजे नैसर्गिक संख्यांची फलने, संगणित करणाऱ्या
ट्यूरिंग यंत्रांत. धन पूर्णांकांचे निरूपण करण्याचा एक सोपा मार्ग म्हणजे
एकाच चिन्हाच्या~$\TMstroke$ क्रमिकांनी त्यांचे सांकेतीकरण करणे.
$n \in \Nat$ असेल, तर $n = 0$ असताना $\TMstroke^n$ ही रिकामी
क्रमिका मानू; अन्यथा ती नेमक्या $n$ $\TMstroke$ चिन्हांची क्रमिका
असेल.
\end{explain}

\begin{defn}[संगणन]
ट्यूरिंग यंत्र~$M$ हे $f\colon \Nat^k \to \Nat$ फलन \emph{संगणित
करते}, असे तेव्हा आणि केवळ तेव्हाच म्हणतो, जेव्हा पुढील स्वरूपाचे
प्रत्येक आदान घेतल्यावर
\[\TMstroke^{n_1} \TMblank \TMstroke^{n_2} \TMblank \dots \TMblank \TMstroke^{n_k}
\]
$M$~थांबते आणि त्याचे प्रदान $\TMstroke^{f(n_1, \dots, n_k)}$ असते.
\end{defn}

\begin{prob}
ट्यूरिंग यंत्र~$M$ हे $f\colon \Nat^k \to \Nat^m$ फलन
कधी संगणित करते याची व्याख्या द्या.
\end{prob}

\begin{ex}\label{tur:mac:una:ex:adder}
\emph{बेरीज:}
$f(n,m) = n + m$ हे फलन संगणित करणारे यंत्र बनवू.
यासाठी सुरुवातीला फितीवर अनुक्रमे $n$ आणि~$m$ लांबीचे
$\TMstroke$ चिन्हांचे दोन खंड असले पाहिजेत आणि शेवटी
$n+m$~$\TMstroke$ चिन्हांचा एकच खंड राहिल्यावर यंत्र थांबले
पाहिजे. दोन $\TMstroke$ आदान-खंडांमध्ये~$\TMblank$ असते. म्हणून एक पद्धत अशी:
$\TMblank$ असलेल्या घरात एक रेघ लिहा आणि शेवटचे~$\TMstroke$
पुसून टाका.
\begin{figure}\[\begin{tikzpicture}[->,>=stealth',shorten >=1pt,auto,node distance=2.8cm,
                    semithick]
  \tikzstyle{every state}=[fill=none,draw=black,text=black]
  \node[initial,state]         (A)              {$q_0$};
  \node[state]         (B) [right of=A] {$q_1$};
  \node[state]         (C) [right of=B] {$q_2$};
  \path (A) edge node {\TMtrans{\TMblank}{\TMstroke}{\TMstay}} (B)
           edge [loop above] node {\TMtrans{\TMstroke}{\TMstroke}{\TMright}} (A)
        (B) edge [loop above] node {\TMtrans{\TMstroke}{\TMstroke}{\TMright}} (B)
            edge node {\TMtrans{\TMblank}{\TMblank}{\TMleft}} (C)
        (C) edge [loop above] node {\TMtrans{\TMstroke}{\TMblank}{\TMstay}} (C);
\end{tikzpicture}
\]\caption{$f(x,y) = x+y$ संगणित करणारे यंत्र}

\label{tur:mac:una:fig:adder}
\end{figure}
\end{ex}

\begin{prob}
$\tuple{3,2}$ हे आदान असताना \ref{tur:mac:una:ex:adder}
मधील यंत्राच्या स्थितिवर्णनांचा क्रम पाहा. यंत्राने $0+0$ संगणित
केल्यास काय घडते?
\end{prob}

\begin{explain}
\ref{tur:mac:rep:ex:doubler} मध्ये आपण एका ट्यूरिंग यंत्राचे उदाहरण दिले:
ते~$\TMstroke$ चिन्हांची क्रमिका आदान म्हणून घेते आणि फितीवर
दुप्पट~$\TMstroke$ चिन्हांची क्रमिका ठेवून थांबते. हे दुप्पट करणारे
यंत्र आहे. पण दुप्पट झालेल्या $\TMstroke$ खंडाच्या डावीकडेही
प्रदानात $\TMblank$ चिन्हे येतात. म्हणून, वाटू शकते तसे, ते
प्रत्यक्षात $f(x) = 2x$ हे फलन संगणित करत नाही. ही अडचण
दूर करण्याचे दोन मार्ग पाहू.
\end{explain}

\begin{ex}
\ref{tur:mac:una:fig:doubler-disc} मधील यंत्र $f(x) =
2x$ हे फलन संगणित करते. \ref{tur:mac:rep:ex:doubler} मधील यंत्र
प्रत्येक आदानातील~$\TMstroke$ पुसून फितीच्या अगदी उजवीकडे दोन
$\TMstroke$ चिन्हे लिहिते. त्याऐवजी हे यंत्र आदानातील प्रत्येक
$\TMstroke$ साठी उजवीकडे एकच~$\TMstroke$ वाढवते. आदान कुठे
संपते हे लक्षात ठेवण्यासाठी ते आदान आणि वाढवलेल्या रेघा यांच्यामध्ये
एक~$\TMblank$ ठेवते; अगदी शेवटी त्या जागी~$\TMstroke$ लिहिते.
तसेच आदानात सध्या कुठे आहोत हे ``लक्षात ठेवण्यासाठी'' आदान-खंडातील
एक~$\TMstroke$ तात्पुरते~$\TMblank$ ने बदलते.
  \begin{figure}
\[\begin{tikzpicture}[->,>=stealth',shorten >=1pt,auto,node distance=2.8cm,
                    semithick]
  \tikzstyle{every state}=[fill=none,draw=black,text=black]
  \node[initial,state] (0)              {$q_0$};
  \node[state]         (1) [above of=0] {$q_1$};
  \node[state]         (2) [above of=1] {$q_2$};
  \node[state]         (3) [right of=2] {$q_3$};
  \node[state]         (4) [below of=3] {$q_4$};
  \node[state]         (5) [below of=4] {$q_5$};
  \node[state]         (6) [above right of=4] {$q_6$};
  \node[state]         (7) [below of=6] {$q_7$};
  \node[state]         (8) [below of=7] {$q_8$};
  \path (0) edge node {\TMtrans{\TMstroke}{\TMblank}{\TMright}} (1)
    (1) edge [loop left] node {\TMtrans{\TMstroke}{\TMstroke}{\TMright}} (1)
      edge node {\TMtrans{\TMblank}{\TMblank}{\TMright}} (2)
    (2) edge [loop above] node {\TMtrans{\TMstroke}{\TMstroke}{\TMright}} (2)
        edge node {\TMtrans{\TMblank}{\TMstroke}{\TMleft}} (3)
    (3) edge [loop above] node {\TMtrans{\TMstroke}{\TMstroke}{\TMleft}} (3)
        edge node[left] {\TMtrans{\TMblank}{\TMblank}{\TMleft}} (4)
    (4) edge        node[left] {\TMtrans{\TMstroke}{\TMstroke}{\TMleft}} (5)
    (5) edge [loop below] node {\TMtrans{\TMstroke}{\TMstroke}{\TMleft}} (5)
        edge              node {\TMtrans{\TMblank}{\TMstroke}{\TMright}} (0)
    (4) edge      node[sloped] {\TMtrans{\TMblank}{\TMstroke}{\TMright}} (6)
    (6) edge              node {\TMtrans{\TMblank}{\TMstroke}{\TMright}} (7)
    (7) edge [loop right] node {\TMtrans{\TMstroke}{\TMstroke}{\TMright}} (7)
        edge              node {\TMtrans{\TMblank}{\TMblank}{\TMleft}} (8)
    (8) edge [loop right] node {\TMtrans{\TMstroke}{\TMblank}{\TMstay}} (8);
    \end{tikzpicture}
\]
\caption{$f(x) = 2x$ संगणित करणारे यंत्र}
\label{tur:mac:una:fig:doubler-disc}
\end{figure}
\end{ex}

\begin{ex}\label{tur:mac:una:ex:mover}
$f(x) = 2x$ संगणित करण्याचा दुसरा मार्ग म्हणजे मूळ दुप्पट
करणारे यंत्र तसेच ठेवून, शेवटी दुप्पट झालेला रेघांचा खंड फितीच्या
अगदी डावीकडे नेणाऱ्या अवस्था आणि सूचना जोडणे. \ref{tur:mac:una:fig:mover}
मधील यंत्र नेमका हा शेवटचा भाग करते: सुरुवातीला फितीवर
प्रथम~$\TMblank$ चिन्हांचा आणि नंतर~$\TMstroke$ चिन्हांचा खंड
असतो (आणि शीर्ष~$\TMblank$ खंडातील कोणत्याही घरावर असते);
मग ते $\TMstroke$ चिन्हे एकेक करून पुसते आणि फितीच्या सुरुवातीला
लिहिते. काम कधी संपले हे ओळखता यावे म्हणून ते प्रथम
$\TMstroke$ खंडाचा शेवट~$\TMendtape$ चिन्हाने दाखवते;
शेवटी हे चिन्ह पुसले जाते. अवस्था~$q_6$ पासून मोजायला
सुरुवात केली आहे, जेणेकरून त्या मूळ दुप्पट करणाऱ्या यंत्रात जोडता
येतील. त्यासाठी एक अतिरिक्त सूचना
$\delta(q_0,
\TMblank) = \tuple{q_6,\TMblank,\TMstay}$, म्हणजे~$q_0$ पासून
~$q_6$ कडे~$\TMtrans{\TMblank}{\TMblank}{\TMstay}$ हे लेबल असलेला
बाण, एवढीच हवी. (येथे एक सूक्ष्म अडचण आहे: परिणामी यंत्र
आदान~$x=0$ साठी चालत नाही. ही अडचण सरावासाठी ठेवू.)
\begin{figure}
  \[  \begin{tikzpicture}[->,>=stealth',shorten >=1pt,auto,node distance=2.8cm,
                      semithick]
    \tikzstyle{every state}=[fill=none,draw=black,text=black]
    \node[initial,state] (6)              {$q_6$};
    \node[state]         (7) [right of=6] {$q_7$};
    \node[state]         (8) [right of=7] {$q_8$};
    \node[state]         (9) [below of=8] {$q_9$};
    \node[state]         (10) [left of=9] {$q_{10}$};
    \node[state]         (11) [left of=10]  {$q_{11}$};
    \node[state]         (12) [below of=11]  {$q_{12}$};
    \node[state]         (13) [right of=12] {$q_{13}$};
    \node[state]         (14) [right of=13] {$q_{14}$};
    \path
    (6)  edge node {\TMtrans{\TMblank}{\TMblank}{\TMright}} (7)
    (7)  edge [loop above] node {\TMtrans{\TMblank}{\TMblank}{\TMright}} (7)
         edge node {\TMtrans{\TMstroke}{\TMstroke}{\TMright}} (8)
    (8)  edge [loop above] node {\TMtrans{\TMstroke}{\TMstroke}{\TMright}} (8)
         edge node {\TMtrans{\TMblank}{\TMendtape}{\TMleft}} (9)
    (9)  edge [loop right] node {\TMtrans{\TMstroke}{\TMstroke}{\TMleft}} (9)
         edge node {\TMtrans{\TMblank}{\TMblank}{\TMright}} (10)
    (10) edge node[above] {\TMtrans{\TMstroke}{\TMblank}{\TMleft}} (11)
    (11) edge [loop above] node {\TMtrans{\TMblank}{\TMblank}{\TMleft}} (11)
         edge node[left] {\begin{tabular}{@{}l@{}}
          \TMtrans{\TMendtape}{\TMendtape}{\TMright}\\
          \TMtrans{\TMstroke}{\TMstroke}{\TMright}
         \end{tabular}} (12)
    (12) edge node[] {\TMtrans{\TMblank}{\TMstroke}{\TMright}} (13)
    (13) edge [loop below] node {\TMtrans{\TMblank}{\TMblank}{\TMright}} (13)
         edge node[sloped] {\TMtrans{\TMstroke}{\TMblank}{\TMleft}} (11)
    (13) edge node {\TMtrans{\TMendtape}{\TMblank}{\TMstay}} (14);
  \end{tikzpicture}
  \]
  \caption{$\TMstroke$ चिन्हांचा खंड डावीकडे नेणे}
  \label{tur:mac:una:fig:mover}
\end{figure}
\end{ex}

\begin{prob}
\ref{tur:mac:una:ex:mover} मध्ये आपण
\ref{tur:mac:rep:ex:doubler} मधील मूळ दुप्पट करणारे यंत्र
आणि \ref{tur:mac:una:fig:mover} मधील खंड हलवणारे यंत्र
एकत्र करून बनणाऱ्या यंत्राचे वर्णन केले. हे संयुक्त यंत्र
आदान~$x=0$, म्हणजे रिकामी फीत, घेऊन सुरू केल्यास काय घडते?
या वेळी यंत्राने प्रदान~$2x=0$ देऊन थांबावे यासाठी त्यात
काय बदल कराल? (एक अवस्था आणि एक संक्रमण जोडून हे करता
आले पाहिजे.)

\end{prob}

\begin{prob}
\emph{वजाबाकी:} $n$ आणि~$m$ लांबीच्या रेघांच्या दोन अरिक्त
क्रमिका आदान म्हणून दिल्या असताना, जेथे $n >
m$ आहे, $f(n,m) = n - m$ हे फलन संगणित करणारे ट्यूरिंग यंत्र
रचा.
\end{prob}

\begin{prob}
\emph{समता:} पुढील फलन संगणित करणारे ट्यूरिंग यंत्र रचा:
\[\fn{equality}(n,m) =
\begin{cases}
  \text{1} & \text{$n = m$ असेल तर} \\
  \text{0} & \text{$n \neq m$ असेल तर}
\end{cases}
\]
येथे~$n$ आणि~$m \in \PosInt$ आहेत.
\end{prob}

\begin{prob}
$x$ आणि $y$ हे धन पूर्णांक फितीवर $\TMblank$ ने विभक्त केलेल्या
$\TMstroke$ चिन्हांच्या क्रमिकांनी निरूपित केले असताना
$\min(x,y)$ हे फलन संगणित करणारे ट्यूरिंग यंत्र रचा.
यंत्राच्या वर्णमालेत अतिरिक्त चिन्हे वापरू शकता.

$\min$ हे फलन आपल्या आदानांतील लहान मूल्य निवडते; म्हणून
$\min(3,5)=3$, $\min(20,16)=16$, $\min(4,4)=4$, इत्यादी.
\end{prob}

\begin{defn}
ट्यूरिंग यंत्र~$M$ हे $f\colon \Nat^k
  \pto \Nat$ हे आंशिक फलन तेव्हा आणि केवळ तेव्हाच संगणित
करते, जेव्हा
  \begin{enumerate}
    \item $f(n_1, \dots, n_k) = m$ असेल तर $M$ हे
    $\TMstroke^{n_1}\concat\TMblank\concat
    \dots \concat\TMblank\concat\TMstroke^{n_k}$ आदानावर
    थांबते आणि त्याचे प्रदान $\TMstroke^{m}$ असते.
    \item $f(n_1, \dots, n_k)$ अपरिभाषित असेल तर $M$
    मुळीच थांबत नाही, किंवा थांबले तरी त्याचे प्रदान
    $\TMstroke$ चिन्हांचा एकच खंड नसते.
  \end{enumerate}
\end{defn}











\section{थांबण्याच्या अवस्था}\label{tur:mac:hal:sec}

\begin{explain}
अमलात आणण्यासाठी कोणतीही सूचना उरली नसेल तेव्हाच आपली यंत्रे
थांबतात, अशी आपण व्याख्या केली आहे. पण ट्यूरिंग यंत्रांच्या
नेहमीच्या निरूपणांत थांबण्यासाठी खास \emph{थांबण्याची अवस्था}~$h$
असते; येथे $h \in Q$.

थांबण्याच्या अवस्थेमागील कल्पना सोपी आहे: यंत्राचे काम संपले की
(ते आदान स्वीकारण्यास तयार झाले किंवा त्याने प्रदान लिहून संपवले
की) ते~$h$ अवस्थेत जाते आणि तिथे थांबते. काही यंत्रांना अशा दोन
अवस्था असतात: एक आदान स्वीकारण्यासाठी आणि दुसरी ते नाकारण्यासाठी.
\end{explain}

\begin{ex}\emph{थांबण्याच्या अवस्था}.
ही कल्पना स्पष्ट करण्यासाठी सम यंत्रात बदल करून पाहू. आदानातील
रेघांची संख्या सम असेल तेव्हा यंत्र~$q_0$ अवस्थेत थांबण्याऐवजी,
त्याला थांबण्याच्या अवस्थेत पाठवणारी सूचना जोडू शकतो.
\[\begin{tikzpicture}[->,>=stealth',shorten >=1pt,auto,node distance=2.8cm,
                    semithick]
  \tikzstyle{every state}=[fill=none,draw=black,text=black]
  \node[initial, state]         (A)                     {$q_0$};
  \node[state]         (B) [right of=A] {$q_1$};
  \node[state]         (C) [below of=A] {$h$};
  \path (A) edge [bend left] node {\TMtrans{\TMstroke}{\TMstroke}{R}} (B)
  	    edge node {\TMtrans{\TMblank}{\TMblank}{N}} (C)
        (B) edge [loop above] node {\TMtrans{\TMblank}{\TMblank}{R}} (B)
            edge [bend left] node {\TMtrans{\TMstroke}{\TMstroke}{R}} (A);
\end{tikzpicture}
\]

उदाहरण आणखी वाढवू. यंत्राला आदानातील रेघांची संख्या विषम असल्याचे
आढळल्यास मूळ सम यंत्र कधीच थांबत नाही. चक्रात फिरत राहणाऱ्या
सूचनेच्या जागी यंत्राला \emph{नकार} अवस्था~$r$ मध्ये नेणारी सूचना
देऊन यंत्रात बदल करू शकतो.
\[\begin{tikzpicture}[->,>=stealth',shorten >=1pt,auto,node distance=2.8cm,
                    semithick]
  \tikzstyle{every state}=[fill=none,draw=black,text=black]
  \node[initial,state]         (A)                     {$q_0$};
  \node[state]         (B) [right of=A] {$q_1$};
  \node[state]         (C) [below of=A] {$h$};
  \node[state]         (D) [below of=B] {$r$};
  \path (A) edge [bend left] node {\TMtrans{\TMstroke}{\TMstroke}{R}} (B)
  	    edge node {\TMtrans{\TMblank}{\TMblank}{N}} (C)
        (B) edge node {\TMtrans{\TMblank}{\TMblank}{N}} (D)
            edge [bend left] node {\TMtrans{\TMstroke}{\TMstroke}{R}} (A);
\end{tikzpicture}
\]
\end{ex}

\begin{explain}
अशा प्रसंगी खास थांबण्याची अवस्था जोडल्याने यंत्र कोणती आदाने
स्वीकारते आणि कोणती नाकारते हे स्पष्ट होते. पण खास थांबण्याची
अवस्था जोडल्याने यंत्राची संगणनशक्ती वाढत नाही, हेही लक्षात
घ्यावे. त्याचप्रमाणे थांबण्याची कमी औपचारिक कल्पनाही उपयोगाची
आहे. या प्रकरणात आतापर्यंत वापरलेल्या थांबण्याच्या व्याख्येमुळे
\emph{थांबण्याची समस्या} याचा पुरावा अंतर्ज्ञानी रीतीने
सहज दाखवता येतो. म्हणून आपण आपली मूळ व्याख्या पुढेही वापरू.
\end{explain}











\section{शिस्तबद्ध यंत्रे}\label{tur:mac:dis:sec}

\begin{explain}
\ref{tur:mac:hal:sec} मध्ये आपण एकच निर्दिष्ट थांबण्याची अवस्था~$h$
असलेली ट्यूरिंग यंत्रे पाहिली: अशी यंत्रे थांबलीच तर~$h$
अवस्थेतच थांबतात. म्हणून एका थांबण्याच्या अवस्थेची यंत्रे, आपण
सर्वसाधारण ट्यूरिंग यंत्रांना जितकी मुभा देतो त्यापेक्षा अधिक
``शिस्तबद्ध'' आहेत. ट्यूरिंग यंत्रांच्या वर्तनावर आणखीही निर्बंध
घालता येतील. उदाहरणार्थ, घर~$0$ वरील फितीच्या डाव्या टोकाचे
चिन्ह पुसण्यास किंवा घर~$0$ वरून डावीकडे जाण्याचा प्रयत्न
करण्यास आपण त्यांना मनाई केलेली नाही. (आपल्या व्याख्येनुसार
अशा वेळी शीर्ष घर~$0$ वरच राहते; इतर व्याख्यांनुसार यंत्र
थांबते.) तसेच यंत्र थांबलेच तर घर~$1$ वर थांबेल, असे गृहीत
धरता आले तर काही वेळा सोयीचे ठरते.
\end{explain}

\begin{defn}\label{tur:mac:dis:defn:disciplined}
ट्यूरिंग यंत्र~$M$ \emph{शिस्तबद्ध} असते तेव्हा आणि केवळ तेव्हाच,
जेव्हा
\begin{enumerate}
    \item त्याला एकच निर्दिष्ट थांबण्याची अवस्था~$h$ असते,
    \item ते थांबलेच तर घर~$1$ वाचत असताना थांबते,
    \item ते घर~$0$ वरील $\TMendtape$ चिन्ह कधीही पुसत नाही, आणि
    \item ते घर~$0$ वरून डावीकडे जाण्याचा कधीही प्रयत्न करत नाही.
\end{enumerate}
\end{defn}

\begin{explain}
कोणतेही ट्यूरिंग यंत्र त्याच वर्तनाचे पण निर्दिष्ट थांबण्याची
अवस्था असलेले यंत्र बनवता येते, हे आपण आधी पाहिले. यासाठी
नवी अवस्था~$h$ जोडायची आणि मूळ~$\delta$ ज्या
$\tuple{q,\sigma}$ जोडीवर अपरिभाषित आहे, त्या प्रत्येक जोडीसाठी
$\delta(q, \sigma) = \tuple{h, \sigma, N}$ ही सूचना जोडायची.
पुरावा किचकट असला तरी कोणतेही ट्यूरिंग यंत्र~$M$ अशा
शिस्तबद्ध ट्यूरिंग यंत्रात~$M'$ बदलता येते की जे त्याच आदानांवर
थांबते आणि तेच प्रदान देते. उदाहरणार्थ, यंत्र थांबत असताना
घर~$1$ वर नसेल, तर शीर्ष फितीच्या डाव्या टोकाचे चिन्ह सापडेपर्यंत
डावीकडे नेणाऱ्या, त्यानंतर एक घर उजवीकडे नेणाऱ्या आणि मग यंत्र
थांबवणाऱ्या काही सूचना जोडता येतात. इतर अटी कशा पूर्ण करता
येतील याचा विचार तुम्ही करा.
मूळ यंत्राने डाव्या टोकाचे चिन्ह पुसले किंवा तेच चिन्ह दुसऱ्या घरावर
लिहिले असेल, तर एवढी साधी शोध-पद्धत पुरेशी नाही; सर्वसाधारण
रूपांतराच्या पुराव्यात हा प्रसंग वेगळा हाताळावा लागतो.

\end{explain}

\begin{ex}
    \ref{tur:mac:dis:fig:adder-disc} मध्ये आपण \ref{tur:mac:una:ex:adder}
    मधील बेरीज-यंत्राचे शिस्तबद्ध यंत्रात रूपांतर करतो.
    \begin{figure}\[\begin{tikzpicture}[->,>=stealth',shorten >=1pt,auto,node distance=2.8cm,
                    semithick]
  \tikzstyle{every state}=[fill=none,draw=black,text=black]
  \node[initial,state]         (A)              {$q_0$};
  \node[state]         (B) [right of=A] {$q_1$};
  \node[state]         (C) [below right of=B] {$q_2$};
  \node[state]         (D) [below left of=C] {$q_3$};
  \node[state]         (H) [left of=D] {$h$};
  \path (A) edge node {\TMtrans{\TMblank}{\TMstroke}{\TMstay}} (B)
           edge [loop below] node {\TMtrans{\TMstroke}{\TMstroke}{\TMright}} (A)
        (B) edge [loop below] node {\TMtrans{\TMstroke}{\TMstroke}{\TMright}} (B)
            edge node {\TMtrans{\TMblank}{\TMblank}{\TMleft}} (C)
        (C) edge node {\TMtrans{\TMstroke}{\TMblank}{\TMleft}} (D)
        (D) edge [loop above] node {\TMtrans{\TMstroke}{\TMstroke}{\TMleft}} (D)
        edge node {\TMtrans{\TMendtape}{\TMendtape}{\TMright}} (H);
\end{tikzpicture}
\]\caption{शिस्तबद्ध बेरीज-यंत्र}

\label{tur:mac:dis:fig:adder-disc}
\end{figure}
\end{ex}

\begin{prop}\label{tur:mac:dis:prop:disciplined} प्रत्येक ट्यूरिंग यंत्र~$M$
साठी असे शिस्तबद्ध ट्यूरिंग यंत्र~$M'$ असते की $M$~प्रदान~$O$
देऊन थांबत असेल तर ते देखील प्रदान~$O$ देऊन थांबते; आणि
$M$~थांबत नसेल तर ते देखील थांबत नाही. विशेषतः, ट्यूरिंग
यंत्राला संगणित करता येणारे कोणतेही $f\colon\Nat^n \to \Nat$
फलन शिस्तबद्ध ट्यूरिंग यंत्रालाही संगणित करता येते.
\end{prop}

\begin{prob}\label{tur:mac:dis:prob:disc-succ}
$f(x) = x+1$ संगणित करणारे शिस्तबद्ध यंत्र द्या.
\end{prob}


\begin{prob}\label{tur:mac:dis:prob:copier}
$\TMstroke^n$ आदानावर सुरू केल्यावर
$\TMstroke^n \concat \TMblank \concat \TMstroke^n$ हे प्रदान
देणारे शिस्तबद्ध यंत्र शोधा.
\end{prob}











\section{ट्यूरिंग यंत्रांचे संयोजन}\label{tur:mac:cmb:sec}

\begin{explain}
आतापर्यंत पाहिलेली ट्यूरिंग यंत्रांची उदाहरणे बरीच सोपी आहेत.
पण आधुनिक कार्यक्रमणभाषेत सोडवता येणारी कोणतीही समस्या ट्यूरिंग
यंत्रांनीही सोडवता येते. अधिक गुंतागुंतीची ट्यूरिंग यंत्रे बनवताना
ती एकत्र जोडता येतात हे समजून घेणे महत्त्वाचे आहे. मग प्रक्रिया
सोप्या भागांत विभागून अधिक गुंतागुंतीच्या समस्या सोडवणारी यंत्रे
बनवता येतात. गुंतागुंतीची समस्या घटक भागांत विभागण्याचा नैसर्गिक
मार्ग सापडला, तर अनेक सोपी ट्यूरिंग यंत्रे तयार करून त्यांचे
एकाच यंत्रात संयोजन करता येते. असे यंत्र अनेक टप्प्यांत समस्या
सोडवते. पुढील विभागातील थांबण्याची समस्या हाताळताना हे विशेष
महत्त्वाचे ठरेल.

$M = \tuple{Q, \Sigma, q_0, \delta}$ आणि
~$M' = \tuple{Q', \Sigma', q_0', \delta'}$ ही ट्यूरिंग यंत्रे
कशी जोडायची? $M$~थांबल्यानंतर त्याच फितीवर
यंत्र~$M'$ ची धाव सुरू करू; फितीवरील मजकूर आणि शीर्षाची जागा कायम राहील.
असे करणारे एकच ट्यूरिंग यंत्र $M \frown M'$ मिळवण्यासाठी पुढील
पायऱ्या घ्या:
\begin{enumerate}
    \item $M'$ च्या~$Q'$ मधील सर्व अवस्थांचे क्रमांक (किंवा नावे)
    बदला, म्हणजे $M$ आणि~$M'$ यांची कोणतीही अवस्था समान राहणार
    नाही ($Q \cap Q' = \emptyset$).
    \item $M \frown M'$ च्या अवस्था $Q \cup Q'$ असतील.
    \item फितीची वर्णमाला $\Sigma \cup \Sigma'$ असेल.
    \item प्रारंभिक अवस्था~$q_0$ असेल.
    \item संक्रमण फलन पुढील $\delta''$ असेल:
    \[\delta''(q,\sigma) =
    \begin{cases}
      \delta(q,\sigma) & \text{$q \in Q$ आणि $\delta(q,\sigma)$ परिभाषित असेल तर}\\
      \delta'(q,\sigma) & \text{$q \in Q'$ असेल तर}\\
      \tuple{q_0', \sigma, \TMstay} & \text{$q \in Q$ असेल आणि
      $\delta(q,\sigma)$ अपरिभाषित असेल तर}
    \end{cases}\]
\end{enumerate}

मिळालेले यंत्र $q \in Q$ अवस्थेत असताना~$M$ च्या सूचना,
तर $q \in Q'$ अवस्थेत असताना~$M'$ च्या सूचना वापरते.
$q \in Q$ अवस्थेत~$\sigma$ चिन्ह वाचताना $M$ मध्ये त्यासाठी
संक्रमण नसेल (म्हणजे $M$ थांबले असते), तर ते~$M'$ च्या
प्रारंभिक अवस्थेत जाते; फितीवरील मजकूर आणि शीर्षाची जागा
तशीच ठेवते.

यंत्र~$M$ शिस्तबद्ध नसेल तर $M$~थांबताना फितीचे शीर्ष कुठे
आहे हे आपल्याला माहीत नसते. म्हणून~$M$ च्या अंतिम
स्थितिवर्णनात शीर्ष घर~$1$ वाचत असेलच असे नाही. यंत्रे
जोडताना हे लक्षात ठेवणे महत्त्वाचे आहे.
\end{explain}

\begin{ex}
\emph{यंत्रांचे संयोजन:} सुरुवातीला $n$ आणि~$m$ लांबीचे
$\TMstroke$ चिन्हांचे दोन खंड आदान म्हणून घेऊन शेवटी
फितीवर $2(m+n)$ $\TMstroke$ चिन्हांचा एकच खंड ठेवून थांबणारे
यंत्र रचू. यासाठी परिचित अशी बेरीज-यंत्र आणि दुप्पट करणारे
यंत्र ही दोन यंत्रे जोडता येतील. आधी बेरीज-यंत्राचा अवस्था-आलेख
काढू.
\[\begin{tikzpicture}[->,>=stealth',shorten >=1pt,auto,node distance=2.8cm,
                    semithick]
  \tikzstyle{every state}=[fill=none,draw=black,text=black]
  \node[initial,state] (A)              {$q_0$};
  \node[state]         (B) [right of=A] {$q_1$};
  \node[state]         (C) [right of=B] {$q_2$};
  \path (A) edge node {\TMtrans{\TMblank}{\TMstroke}{\TMstay}} (B)
            edge [loop above] node {\TMtrans{\TMstroke}{\TMstroke}{\TMright}} (A)
        (B) edge [loop above] node {\TMtrans{\TMstroke}{\TMstroke}{\TMright}} (B)
            edge node {\TMtrans{\TMblank}{\TMblank}{\TMleft}} (C)
        (C) edge [loop above] node {\TMtrans{\TMstroke}{\TMblank}{\TMstay}} (C);
\end{tikzpicture}
\]
प्रदान दुप्पट करण्यासाठी अवस्था~$q_2$ मध्ये थांबण्याऐवजी
यंत्र पुढे चालवायचे आहे. दुप्पट करणारे यंत्र आदानातील पहिली
रेघ पुसते आणि स्वतंत्र प्रदानात दोन रेघा लिहिते, हे आठवा.
बेरीज-यंत्राच्या प्रदानातील पहिली रेघ शीर्ष वाचत असेल याची
खात्री करणारी सूचना जोडू.
\[\begin{tikzpicture}[->,>=stealth',shorten >=1pt,auto,node distance=2.8cm,
                    semithick]
  \tikzstyle{every state}=[fill=none,draw=black,text=black]
  \node[initial,state] (A)              {$q_0$};
  \node[state]         (B) [right of=A] {$q_1$};
  \node[state]         (C) [right of=B] {$q_2$};
  \node[state]         (D) [below left of=C] {$q_3$};
  \node[state]         (E) [below left of=D] {$q_4$};
  \path (A) edge node {\TMtrans{\TMblank}{\TMstroke}{\TMstay}} (B)
            edge [loop above] node {\TMtrans{\TMstroke}{\TMstroke}{\TMright}} (A)
        (B) edge [loop above] node {\TMtrans{\TMstroke}{\TMstroke}{\TMright}} (B)
            edge node {\TMtrans{\TMblank}{\TMblank}{\TMleft}} (C)
        (C) edge node {\TMtrans{\TMstroke}{\TMblank}{\TMleft}} (D)
        (D) edge [loop left] node {\TMtrans{\TMstroke}{\TMstroke}{\TMleft}} (D)
            edge node[left, xshift=-2mm] {\TMtrans{\TMendtape}{\TMendtape}{\TMright}} (E);
\end{tikzpicture}
\]
आता आदान दुप्पट करणे सोपे आहे: दुप्पट करणारे यंत्र
अवस्था~$q_4$ ला जोडायचे. त्यासाठी त्या यंत्राच्या अवस्थांना
पुन्हा नावे द्यावी लागतील, म्हणजे त्या~$q_0$ ऐवजी~$q_4$
पासून सुरू होतील. अशा रीतीने दोन प्रारंभिक अवस्था राहणार नाहीत.
अंतिम अवस्था-आलेख \ref{tur:mac:cmb:fig:combined} मध्ये दाखवल्याप्रमाणे
असेल.
\begin{figure}
\[\begin{tikzpicture}[->,>=stealth',shorten >=1pt,auto,node distance=2.8cm,
                    semithick]
  \tikzstyle{every state}=[fill=none,draw=black,text=black]
  \node[initial,state] (A)              {$q_0$};
  \node[state]         (B) [right of=A] {$q_1$};
  \node[state]         (C) [right of=B] {$q_2$};
  \node[state]         (D) [below left of=C] {$q_3$};
  \node[state]         (E) [below left of=D] {$q_4$};
  \node[state]         (2) [right of=E] {$q_5$};
  \node[state]         (3) [right of=2] {$q_6$};
  \node[state]         (4) [below of=3] {$q_7$};
  \node[state]         (5) [left of=4]  {$q_8$};
  \node[state]         (6) [left of=5]  {$q_9$};
  \path (A) edge node {\TMtrans{\TMblank}{\TMstroke}{\TMstay}} (B)
            edge [loop above] node {\TMtrans{\TMstroke}{\TMstroke}{\TMright}} (A)
        (B) edge [loop above] node {\TMtrans{\TMstroke}{\TMstroke}{\TMright}} (B)
            edge node {\TMtrans{\TMblank}{\TMblank}{\TMleft}} (C)
        (C) edge node {\TMtrans{\TMstroke}{\TMblank}{\TMleft}} (D)
        (D) edge [loop left] node {\TMtrans{\TMstroke}{\TMstroke}{\TMleft}} (D)
            edge node[left, xshift=-2mm] {\TMtrans{\TMendtape}{\TMendtape}{\TMright}} (E)
    (E) edge node {\TMtrans{\TMstroke}{\TMblank}{\TMright}} (2)
    (2) edge [loop above] node {\TMtrans{\TMstroke}{\TMstroke}{\TMright}} (2)
      edge node {\TMtrans{\TMblank}{\TMblank}{\TMright}} (3)
    (3) edge [loop above] node {\TMtrans{\TMstroke}{\TMstroke}{\TMright}} (3)
        edge node {\TMtrans{\TMblank}{\TMstroke}{\TMright}} (4)
    (4) edge [loop below] node {\TMtrans{\TMblank}{\TMstroke}{\TMleft}} (4)
        edge node {\TMtrans{\TMstroke}{\TMstroke}{\TMleft}} (5)
    (5) edge [loop below]  node {\TMtrans{\TMstroke}{\TMstroke}{\TMleft}} (5)
        edge              node {\TMtrans{\TMblank}{\TMblank}{\TMleft}} (6)
    (6) edge [loop below] node {\TMtrans{\TMstroke}{\TMstroke}{\TMleft}} (6)
        edge              node {\TMtrans{\TMblank}{\TMblank}{\TMright}} (E);
\end{tikzpicture}
\]
\caption{बेरीज-यंत्र आणि दुप्पट करणारे यंत्र यांचे संयोजन}

\label{tur:mac:cmb:fig:combined}
\end{figure}
\end{ex}

\begin{prop}
    $M$ आणि $M'$ ही शिस्तबद्ध यंत्रे अनुक्रमे $f\colon
    \Nat^k \to \Nat$ आणि $f'\colon \Nat \to \Nat$ ही फलने
    संगणित करत असतील, तर $M \frown M'$ हेही शिस्तबद्ध असून
    ~$\comp{f}{f'}$ संगणित करते.
\end{prop}

\begin{proof}
    $M$ शिस्तबद्ध असल्यामुळे ते
    प्रदान~$f(n_1,\dots,n_k) = m$ देऊन थांबते तेव्हा शीर्ष
    घर~$1$ वाचत असते. आता यंत्र~$M'$ च्या प्रारंभिक अवस्थेत
    गेल्यावर प्रदान $f'(m)$ देऊन यंत्र पुन्हा घर~$1$ वरच थांबते.
    \ref{tur:mac:dis:defn:disciplined} मधील इतर अटीही पूर्ण होतात.
\end{proof}

\begin{prob}
    \ref{tur:mac:dis:prob:disc-succ} मधील यंत्र~$M$ घेऊन
    $M \frown M$ रचा आणि $f(x) = x+2$ संगणित करणारे
    शिस्तबद्ध ट्यूरिंग यंत्र द्या.
\end{prob}










\section{ट्यूरिंग यंत्रांचे विविध प्रकार}\label{tur:mac:var:sec}

खरे तर ट्यूरिंग यंत्रांची व्याख्या अनेक प्रकारे करता येते;
आपली व्याख्या हा त्यांपैकी एकच प्रकार आहे. काही बाबतींत ती
इतर व्याख्यांपेक्षा अधिक मोकळी आहे. आपण कोणतीही सांत वर्णमाला
चालू देतो; अधिक मर्यादित व्याख्येत फितीवरील फक्त दोन चिन्हे,
$\TMstroke$ आणि~$\TMblank$, चालतील. आपण यंत्राला एकाच वेळी
फितीवर चिन्ह लिहू आणि शीर्ष हलवू देतो; इतर व्याख्यांत लिहिणे
किंवा शीर्ष हलवणे यांपैकी एकच क्रिया करता येते. शीर्ष न सरकवता
लिहिण्याची शक्यताही आपण ठेवतो; इतर व्याख्यांत $\TMstay$
ही ``सूचना'' नसते. दुसरीकडे, काही बाबतींत आपली व्याख्या अधिक
मर्यादित आहे. फीत फक्त एका दिशेला अनंत आहे असे आपण मानले;
इतर व्याख्यांत ती डावीकडे आणि उजवीकडे दोन्ही दिशांना अनंत
असते. एवढेच नव्हे, तर कितीही सांत संख्येने स्वतंत्र फिती किंवा घरांची
अनंत जाळीही चालू देता येईल. आपण ट्यूरिंग यंत्राचा सूचनासंच
संक्रमण फलनाने दर्शवतो; इतर व्याख्यांत संक्रमण संबंध वापरतात,
ज्यामुळे एखाद्या परिस्थितीत यंत्रासाठी एकाहून अधिक सूचना
शक्य असतात.

व्याख्येतील हा शेवटचा बदल विशेष लक्षवेधक आहे. आपल्या व्याख्येत
यंत्र~$q$ अवस्थेत~$\sigma$ चिन्ह वाचत असताना
$\delta(q, \sigma)$ पुढील चिन्ह, अवस्था आणि शीर्षाची जागा
निश्चित करते. पण सूचनासंच हा सध्याच्या अवस्था-चिन्ह जोड्या
$\tuple{q,
  \sigma}$ आणि पुढील अवस्था-चिन्ह-दिशा त्रयी $\tuple{q', \sigma',
  D}$ यांमधील संबंध म्हणून घेतला, तर यंत्राची कृती एकमेवपणे
निश्चित होणार नाही: सूचना-संबंधात $\tuple{q,
  \sigma, q', \sigma', D}$ आणि $\tuple{q, \sigma, q'', \sigma'', D'}$
या दोन्ही नोंदी असू शकतात. अशा वेळी आपल्याला
\emph{अनिर्धारक} ट्यूरिंग यंत्र मिळते. संगणकीय गुंतागुंत
सिद्धांतात अशा यंत्रांना महत्त्वाचे स्थान आहे.

ट्यूरिंग यंत्र कधी थांबते याबाबतही वेगवेगळे संकेत आहेत:
संक्रमण फलन अपरिभाषित असेल तेव्हा ते थांबते असे आपण मानतो;
इतर व्याख्यांत यंत्राने खास निर्दिष्ट थांबण्याच्या अवस्थेत
असणे आवश्यक असते. अशी निर्दिष्ट अवस्था मागितल्याने ट्यूरिंग
यंत्रांना काय संगणित करता येते यावर निर्बंध येत नाही, हे आपण
\ref{tur:mac:hal:sec} मध्ये स्पष्ट केले आहे. आपल्या यंत्रांची फीत
एकाच दिशेला अनंत असल्यामुळे कधी एखादी सूचना नीट अमलात
आणता येत नाही: यंत्र सर्वांत डावीकडील घर वाचत आहे आणि
डावीकडे सरकायची सूचना आहे. आपल्या व्याख्येनुसार शीर्ष
फितीबाहेर ``पडण्याऐवजी'' त्याच घरावर राहते; पण त्या वेळी
यंत्र थांबेल अशीही व्याख्या करता आली असती. ती व्याख्याही
सममूल्य आहे: घर~$0$ वरून डावीकडे जाण्याचा प्रयत्न केल्यावर
थांबणाऱ्या यंत्राचे वर्तन अनुकरण करण्यासाठी प्रत्येक
$\delta(q,\TMendtape) =
\tuple{q',\sigma,\TMleft}$ संक्रमण काढून टाकता येईल.
मग~$\TMendtape$ वरून डावीकडे जाण्याचा प्रयत्न करण्याऐवजी
यंत्र थांबेल.\footnote{ही पद्धत \emph{अगदी तशी} चालत नाही,
कारण घर~$0$ सोडून इतर घरांतही $\TMendtape$ लिहिण्यास किंवा
वाचण्यास मनाई नाही (\ref{tur:mac:una:ex:mover} पाहा). अशा वापरासाठी
दुसरे~$\TMendtape'$ चिन्ह जोडून ही अडचण दूर करता येते.
तसेच मूळ यंत्राने घर शून्यवरील चिन्ह बदलले असल्यास त्या सीमास्थानाची
ओळख जतन करण्यासाठी स्वतंत्र स्थिर सीमा-चिन्ह किंवा योग्य सांकेतिक
व्यवस्था आवश्यक आहे; फक्त इतर घरांसाठी दुसरे चिन्ह वापरणे पुरेसे नाही.}


संख्या निरूपित करण्याचेही वेगवेगळे मार्ग आहेत (आणि त्यामुळे
ट्यूरिंग यंत्र संगणित करत असलेल्या आदान-प्रदान फलनाचे रूपही
बदलते): आपण एकचिन्ही निरूपण वापरतो, पण द्विमान निरूपणही
वापरता येते. त्यासाठी $\TMblank$ आणि~$\TMendtape$ यांच्या
जोडीला आणखी दोन चिन्हे लागतात.

आता एक महत्त्वाची गोष्ट: यांपैकी कोणत्याही बदलामुळे नेमकी
कोणती फलने ट्यूरिंग-संगणनक्षम आहेत हे बदलत नाही.
\emph{एका व्याख्येनुसार एखादे फलन ट्यूरिंग-संगणनक्षम असेल,
तर त्या सर्व व्याख्यांनुसार ते ट्यूरिंग-संगणनक्षम असते.}

याची सविस्तर पडताळणी येथे करणार नाही. एक उदाहरण पाहू:
दोन दिशांना अनंत असलेली फीत किंवा अनेक फिती वापरू दिल्याने
संगणनशक्ती वाढत नाही. कारण साधारणपणे एकाच दिशेला अनंत असलेली
एक फीत वापरणारे ट्यूरिंग यंत्र अनेक फितींचे किंवा दोन्ही
दिशांना अनंत असलेल्या फितींचे अनुकरण करू शकते. उदाहरणार्थ,
अतिरिक्त अवस्था आणि सूचना वापरून अनेक फितींच्या किंवा
दोन दिशांना अनंत असलेल्या फितीच्या यंत्राचा कार्यक्रम, एका
दिशेला अनंत असलेली एकच फीत वापरणाऱ्या यंत्राच्या कार्यक्रमात
``रूपांतरित'' करता येतो. रूपांतरित यंत्र सम क्रमांकांची घरे
फीत~$1$ वरील घरांसाठी (किंवा द्विदिश अनंत फितीच्या ``धन''
बाजूच्या घरांसाठी), आणि विषम क्रमांकांची घरे फीत~$2$ वरील
घरांसाठी (किंवा त्या फितीच्या ``ऋण'' बाजूच्या घरांसाठी)
वापरू शकते.











\section{चर्च--ट्यूरिंग प्रबंध}\label{tur:mac:ctt:sec}

प्रभावी प्रक्रियेच्या संकल्पनेला अचूक पर्याय म्हणून ट्यूरिंग
यंत्रे मांडली जातात. ट्यूरिंग यांच्या मते, प्रभावी प्रक्रिया आणि
ट्यूरिंग यंत्र या दोन्ही संकल्पना समजलेल्या कोणालाही प्रभावी
प्रक्रियेने करता येणारी कोणतीही गोष्ट ट्यूरिंग यंत्रानेही करता
येईल असे अंतर्ज्ञान होईल. प्रभावी प्रक्रियेच्या संकल्पनेचे
इतर सर्व प्रस्तावित अचूक पर्याय व्याप्तीच्या दृष्टीने ट्यूरिंग
यंत्राच्या संकल्पनेशी सममूल्य ठरतात, ही गोष्ट या दाव्याला
आधार देते---म्हणजे ते नेमका तोच फलनांचा संच संगणित करू
शकतात. या दाव्याला \emph{चर्च--ट्यूरिंग प्रबंध} म्हणतात.

\begin{defn}[चर्च--ट्यूरिंग प्रबंध]
\emph{चर्च--ट्यूरिंग प्रबंध} असे सांगतो की प्रभावी प्रक्रियेने
संगणित करता येणारी कोणतीही गोष्ट ट्यूरिंग-संगणनक्षम असते.
\end{defn}

चर्च--ट्यूरिंग प्रबंधाचा दोन प्रकारे आधार घेतला जातो. पहिला
वापर काहीसा आळशीपणाला सबब देणारा आहे. समजा, एखादी गोष्ट
संगणित करण्याच्या प्रभावी प्रक्रियेचे वर्णन आपल्याकडे
``छद्म-संकेतलेखनात'' आहे. अशा वेळी प्रत्यक्षात ट्यूरिंग यंत्र
बनवले नसले तरी हेच फलन एखादे ट्यूरिंग यंत्र संगणित करते,
असा दावा चर्च--ट्यूरिंग प्रबंधाच्या आधारे समर्थित करता येतो.

चर्च--ट्यूरिंग प्रबंधाचा दुसरा वापर तत्त्वज्ञानाच्या दृष्टीने
अधिक रोचक आहे. काही फलने ट्यूरिंग यंत्रांनी संगणित करता
येत नाहीत, हे दाखवता येते. त्यावरून चर्च--ट्यूरिंग प्रबंध
वापरून कोणत्याही प्रक्रियेने ती प्रभावी रीतीने संगणित करता
येत नाहीत, असा निष्कर्ष काढता येतो. कारण तशी एखादी प्रक्रिया
असती तर चर्च--ट्यूरिंग प्रबंधानुसार त्यासाठी ट्यूरिंग यंत्रही
असले असते. म्हणून एखादे फलन संगणित करणारे ट्यूरिंग यंत्र
असू शकत नाही हे सिद्ध केले, तर त्यासाठी प्रभावी प्रक्रियाही
असू शकत नाही. विशेषतः, तथाकथित थांबण्याची समस्या केवळ
ट्यूरिंग यंत्रांनीच सोडवता येत नाही असे नव्हे; ती मुळीच
प्रभावी रीतीने सोडवता येत नाही, असा दावा करण्यासाठी
चर्च--ट्यूरिंग प्रबंध वापरला जातो.



